
\section{Worked module patterns}
\label{sec:examples}

\subsection{Certified rational bisection for \texorpdfstring{\(y^3+y=q\)}{y3+y=q}}

This example shows exact arithmetic, branching, iteration, and a quantitative
output without invoking a completed real continuum.

Let
\[
     F(y)=y^3+y-q
\]
for rational \(q\).  Suppose rational endpoints \(L_0<U_0\) satisfy
\[
       F(L_0)\le0\le F(U_0).
\]
All rational comparisons are decidable by integer cross multiplication.

Define
\[
  M_t=\frac{L_t+U_t}{2}
\]
and update
\[
 (L_{t+1},U_{t+1})=
 \begin{cases}
  (M_t,U_t),&F(M_t)\le0,\\
  (L_t,M_t),&F(M_t)>0.
 \end{cases}
\]
The state is the pair of rational endpoints.  The branch bit is the reflected
QF inequality \(F(M_t)\le0\).

\begin{lemma}[Strict monotonicity certificate]
For rationals \(u<v\),
\[
      u^3+u<v^3+v.
\]
\end{lemma}

\begin{proof}
The exact identity
\[
 (v^3+v)-(u^3+u)
   =(v-u)(v^2+uv+u^2+1)
\]
is a ring calculation.  The first factor is positive.  For the second,
\[
 v^2+uv+u^2
   = (v-u)^2+3uv
\]
is not a uniformly useful positivity decomposition when \(uv<0\), but
\[
 4(v^2+uv+u^2)
    =(2v+u)^2+3u^2
\]
is nonnegative.  Hence \(v^2+uv+u^2+1>0\).  Integer cross multiplication
internalizes the rational inequalities.
\end{proof}

Use the quantifier-free invariant
\[
\Inv(L,U,t):
  L<U\ \wedge\
  F(L)\le0\ \wedge\
  0\le F(U)\ \wedge\
  2^t(U-L)=U_0-L_0.
\]
The last equation may be replaced by an inequality if denominator
normalization makes exact equality inconvenient.

\begin{theorem}[Bisection contract]
For every \(T\), the run returns \(L_T,U_T\) satisfying the invariant and
\[
       U_T-L_T=2^{-T}(U_0-L_0).
\]
Moreover there is at most one rational \(y\in[L_T,U_T]\) with \(F(y)=0\).
\end{theorem}

\begin{proof}
The base is the input certificate.  For the step, the midpoint lies strictly
between the endpoints.  In the first branch, \(F(M)\le0\), so the new lower
endpoint preserves the bracket; in the second branch, decidability yields
\(F(M)>0\), so the new upper endpoint preserves it.  The width halves by
rational arithmetic.  Apply fixed trace reflection.

If \(u<v\) were two roots in the interval, strict monotonicity would give
\(F(u)<F(v)\), contradicting both being zero.
\end{proof}

The midpoint \(a_T=(L_T+U_T)/2\) is a finite observation.  Choosing
\[
 T(k)=k+\left\lceil\log_2(U_0-L_0)\right\rceil+1
\]
gives an interval of width at most \(2^{-k}\).  The sequence \(a_{T(k)}\) is a
regular rational name, and its modulus is explicit.  Existence of an abstract
real root is not needed to prove every finite enclosure; alternatively, the
name itself may be taken as the constructive root object.

\paragraph{Lean path.}
The implementation uses a rational setoid, a state packed by \code{pairT}, a
step term built from a Boolean inequality bit, \code{recDef} for the run,
\code{ipra_ring} for the width identity, and \code{IPRA.ind} for the invariant.
No general real-analysis library is required.

\subsection{Complete bounded rational-point enumeration}

Let a curve or variety code \(C\) determine a decidable polynomial equation
\(W(C,x)=1\) on canonical rational tuples.  Let \(h(C,x)\) be a natural height.

Assume modules:
\[
\begin{aligned}
 \mathsf{height}:&
 \quad W(C,x)=1\vdash h(C,x)\le H(C),\\
 \mathsf{cover}:&
 \quad h(C,x)\le H(C)\vdash
   \exists j<N(C)\ [E(C,j)\sim x],
\end{aligned}
\]
where \(E\) enumerates canonical bounded tuples and \(\sim\) is code
equivalence.

The consumer scans \(j<N(C)\), checks \(W(C,E(C,j))\), and appends qualifying
canonical codes to a duplicate-free list.

\begin{theorem}[Bounded point-list correctness]
From \(\mathsf{height}\) and \(\mathsf{cover}\), iPRA proves:
\[
 W(C,x)=1
 \quad\iffc\quad
 \exists j<\ell(L_C)\ [L_C[j]\sim x],
\]
where \(L_C\) is the computed finite list.
\end{theorem}

\begin{proof}
Soundness follows from the append condition and extensionality of \(W\).
For completeness, the height module gives the bounded region and the cover
module gives an index \(j<N(C)\).  A scan invariant says:
\[
 \forall i<r,\quad
 W(C,E(C,i))=1
 \imp E(C,i)\text{ is represented in the partial list}.
\]
Compile the bounded universal statement to a Boolean product.  At the next
scan step, either the candidate fails \(W\), or it is already equivalent to an
entry, or it is appended.  All cases preserve the invariant.  At \(r=N(C)\),
the covering index is included.
\end{proof}

The consumer is complete and object-theoretic.  If the height module is
unknown, the theorem remains
\[
  \mathsf{HeightBound}(C,H),\mathsf{Cover}(C,H)
   \vdash_{\PRAAx}\mathsf{CompleteList}(C,L_C).
\]
This is the exact form appropriate for general Faltings- and Siegel-related
packets.

\subsection{A finite nerve and certified homology observation}

Let \(X\) be represented at scale \(k\) by centers
\(x_0,\ldots,x_{m-1}\) and rational radii.  A certificate decides for every
subset \(\sigma\subseteq\{0,\ldots,m-1\}\) of size at most \(d+1\) whether the
corresponding finite intersection is declared nonempty.  Assume the declarations
satisfy downward closure.

Encode a simplex by a bit mask \(s<2^m\).  The predicate
\[
   \mathsf{Simplex}(s)=1
\]
checks the mask size and the supplied intersection certificate.  Boundary
matrices over \(\mathbb F_2\) are bounded arrays.  Gaussian elimination is a
fixed finite trace with an invariant that every row operation preserves the
kernel and image dimensions.

\begin{theorem}[Finite nerve homology]
PCFM derives the ranks of the boundary maps, verifies
\(\partial_{r-1}\partial_r=0\), and computes
\[
   \beta_r
     = \dim C_r-\operatorname{rank}\partial_r
       -\operatorname{rank}\partial_{r+1}
\]
for the certified finite nerve.
\end{theorem}

\begin{proof}
The chain-complex equation is a bounded parity count: every codimension-two
face occurs twice.  Boolean finite folds verify each matrix entry.  Row
reduction uses exact XOR row operations; the trace invariant states that the
current matrix is related to the original by invertible row operations and
that pivot columns are distinct.  At terminal form, bounded search counts
pivots.  The displayed dimension formula is integer arithmetic.
\end{proof}

To identify this finite answer with a topological invariant of \(X\), one needs
a producer such as a good-cover theorem, a reach/separation bound, or a
certified interleaving showing stability over a scale interval.  That producer
is not part of finite linear algebra and remains explicit.

\subsection{A truncated hypergeometric identity}

Consider the finite polynomial
\[
  S_N(a,b;c;z)
   =\sum_{n<N}
      \frac{(a)_n(b)_n}{(c)_n\,n!}z^n
\]
under rational parameter guards that keep every denominator nonzero.
Pochhammer symbols, powers, factorials, and the finite sum are primitive
recursive rational terms.

A contiguous relation has the form
\[
  A(a,b,c,z)S_N(\theta_1)
  +B(a,b,c,z)S_N(\theta_2)
  +C(a,b,c,z)S_N(\theta_3)
  =R_N,
\]
where \(R_N\) is a boundary term.  After multiplying by a common denominator,
the statement is an integer polynomial identity.

\begin{theorem}[Finite contiguous relation]
For every \(N\), the exact truncated identity with boundary term \(R_N\) is
derivable in iPRA.
\end{theorem}

\begin{proof}
Prove the equality at \(N=0\).  At the successor step subtract the \(N\)-case
from the \(N+1\)-case.  Only the newly added summands and the update of
\(R_N\) remain.  Clear denominators using the guard, normalize the resulting
polynomial identity, and apply quantifier-free induction on \(N\).
\end{proof}

To obtain the infinite analytic identity at precision \(2^{-k}\), add tail
modules for the three series and choose \(N(k)\) so that the weighted sum of
tail bounds is at most \(2^{-k}\).  This cleanly separates the universally
reusable finite algebra from domain-specific analytic estimates.

\subsection{Guarded Euler--Rodrigues axis extraction}

For a rational or regular-name rotation matrix \(R\), the antisymmetric part
suggests an axis vector
\[
   v=(R_{32}-R_{23},R_{13}-R_{31},R_{21}-R_{12}).
\]
At the identity rotation, \(v=0\) and no oriented axis is determined.  An
unrestricted exact-axis selector is therefore constructively invalid.

PCFM offers two correct endpoints.

\paragraph{Guarded exact endpoint.}
Assume a separation certificate \(\|v\|\ge 2^{-g}\).  Normalize \(v\) using a
guarded reciprocal/square-root module and prove the resulting vector is an
eigenvector with eigenvalue \(1\), up to requested precision.

\paragraph{Stable relational endpoint.}
Return the projective line or a finite antipodal pair of approximate axes.
Observations invariant under sign, such as the rank-one projector
\(vv^\top/\|v\|^2\), are quotient-stable.  Near the identity return the
rotation matrix itself or an error ball rather than choosing an axis.

This example illustrates the framework's policy: repair an invalid global
selection theorem by changing the output interface, not by hiding a classical
choice inside a program.
